\documentclass[10pt,a4paper]{amsart}
\usepackage[margin=19mm]{geometry}
\usepackage{amsmath,amssymb,mathtools}
\usepackage[scr=rsfs,cal=euler]{mathalfa}
\usepackage{xfrac}
\usepackage[unicode,hidelinks,bookmarks=false]{hyperref}
\newcommand{\ie}{i.e.\ }
\newcommand{\cf}{cf.\ }
\newcommand{\resp}{respectively\ }
\newcommand{\Subsection}{\subsection}
\providecommand{\nobreakdash}{\nobreak\hbox{-}}
\setlength{\emergencystretch}{2em}
\multlinegap0pt
\usepackage{bm}
\usepackage{bbm}
\usepackage{dsfont}
\usepackage{xspace}
\usepackage{cancel}
\usepackage{mathtools}
\usepackage{amsthm}
\makeatletter
\@ifundefined{theorem}{\newtheorem{theorem}{Theorem}}{}
\@ifundefined{proposition}{\newtheorem{proposition}[theorem]{Proposition}}{}
\makeatother
\newcommand{\CC}{\mathcal{C}}
\newcommand{\HH}{\mathcal{H}}
\newcommand{\di}{\mathsf{d}}
\newcommand{\maxi}{\mathsf{max}}
\title[Countable dense homogeneity in large products of Polish spac]{Countable dense homogeneity in large products of Polish spaces}
\author{Andrea Medini, Juris Stepr\=ans}
\date{Source: arXiv:2510.25322v1 (CC BY 4.0)}
\begin{document}
\maketitle

\noindent\emph{Source and licence.} Countable dense homogeneity in large products of Polish spaces. Version arXiv:2510.25322v1 (CC BY 4.0), distributed under the Creative Commons Attribution 4.0 International licence, as recorded in the arXiv licence metadata for this exact version and shown on the abstract page https://arxiv.org/abs/2510.25322v1. Full rights notice: licenses/arxiv-2510.25322-CC-BY-4.0.txt.\par\smallskip

\noindent\textit{Editorial scope.} Counted unit: the Theorem whose source label is \texttt{theorem\_convergence} in the article (source0.tex lines 226--233, source label \texttt{theorem\_convergence}), delivered together with the article's own complete original proof of it (source0.tex lines 234--282, one \texttt{proof} environment of 49 lines). No mathematics is written, repaired or re-derived: statement and proof are the article's own text line for line. Required context retained uncounted: proposition\_convergence (lines 179--181). Every definition, display and label of those lines is retained unchanged. Omitted from this package and not redistributed: 1--178, 182--225, 283--989. Nothing inside a retained excerpt is omitted or altered: every retained body is delivered exactly as the article's lines are written, so no omission receipt of any kind accompanies this package. Citations: the retained excerpts carry no citation command at all, so no bibliography entry of the article is transcribed and no reference is redistributed. Reference adaptation: none was needed and none was made; every reference inside the delivered unit resolves inside it, so no unresolved reference or citation is produced. Printed slips kept literal: no printed word, symbol, hypothesis or number of the counted statement or of its proof was altered. Layout and class: the article's own class is \texttt{amsart} with packages including amssymb, amsmath, mathrsfs, verbatim, url, bbm, bbold; the wrapper is the lane's pinned \texttt{amsart} editorial wrapper at 10pt on A4, which restates the theorem-like environments the retained text needs and adds the article's own macro definitions that the retained text needs (\texttt{\textbackslash CC}, \texttt{\textbackslash HH}, \texttt{\textbackslash di}, \texttt{\textbackslash maxi}), with extra wrapper packages (bm, bbm, dsfont, xspace, cancel, mathtools). The delivered numbering is the unit's own. Assets: no retained span contains a figure environment, a graphics-inclusion command, a PDF-page inclusion, a table or a TikZ picture; no image file is copied into this package, the package declares no asset dependency and no image file is redistributed with it. Licence: version arXiv:2510.25322v1 of this article is distributed under the Creative Commons Attribution 4.0 International licence, as recorded in the arXiv licence metadata for this exact version and shown on the abstract page https://arxiv.org/abs/2510.25322v1. A full rights notice is delivered at licenses/arxiv-2510.25322-CC-BY-4.0.txt. Characters: every retained excerpt is pure ASCII, so the delivered engine, XeLaTeX with its default Latin Modern text font, reports no missing character.
\begin{proposition}\label{proposition_convergence}
Let $X$ be a space, let $Y$ be a complete metric space with distance~$\di$, and let $f_n\in\CC(X,Y)$ for $n\in\omega$. Assume that $\di\big(f_{n+1}(x),f_n(x)\big)\leq 2^{-n}$ for every $n\in\omega$ and $x\in X$. Then $(f_n:n\in\omega)$ converges uniformly (to a continuous function).
\end{proposition}

\begin{theorem}\label{theorem_convergence}
Let $X$ be a complete metric space with distance $\di$, and let $h_n\in\HH(X)$ for $n\in\omega$. Set $H_n=h_n\circ\cdots\circ h_0$ for $n\in\omega$. Assume that the following conditions hold for each $n$:
\begin{enumerate}
\item $\di\big(h_{n+1}(x),x\big)\leq 2^{-n}$ for every $x\in X$,
\item $\di\big(H_{n+1}^{-1}(x),H_n^{-1}(x)\big)\leq 2^{-n}$ for every $x\in X$.
\end{enumerate}
Then $(H_n:n\in\omega)$ converges uniformly to an element of $\HH(X)$.
\end{theorem}

\begin{proof}
Proposition \ref{proposition_convergence} and condition $(1)$ ensure that $(H_n:n\in\omega)$ converges uniformly to some $H\in\CC(X,X)$. Similary, Proposition \ref{proposition_convergence} and condition $(2)$ ensure that $(H^{-1}_n:n\in\omega)$ converges uniformly to some $H^\ast\in\CC(X,X)$. It remains to show that $H^\ast$ is the inverse of $H$. So fix $x\in X$, then set
$$
x_{m,n}=H_m^{-1}\big(H_n(x)\big)\text{ and }y_{m,n}=H_m\big(H_n^{-1}(x)\big).
$$
for $m,n\in\omega$. Clearly, the following four claims will conclude the proof.

\noindent\textbf{Claim 1.} $x_{m,n}\to x$.

\noindent\textit{Proof.} Pick $\varepsilon >0$. Fix $K\in\omega$ such that $2^{-K}+2^{-(K+1)}+\cdots <\varepsilon$. Assume that $m>n\geq K$, and set $y=H_n(x)$. Then
\begin{align}
\di(x_{m,n},x)&= \di\big(H_m^{-1}(y),H_n^{-1}(y)\big)\nonumber\\
&\leq\di\big(H_m^{-1}(y),H_{m-1}^{-1}(y)\big)+\cdots +\di\big(H_{n+1}^{-1}(y),H_n^{-1}(y)\big)\nonumber\\
&\leq 2^{-(m-1)}+\cdots +2^{-n}\nonumber\\
&<\varepsilon\nonumber
\end{align}
by condition $(2)$. A similar argument yields the same inequality when $n>m\geq K$, and the case $m=n\geq K$ is trivial. $\blacksquare$

\noindent\textbf{Claim 2.} $x_{m,n}\to H^\ast\big(H(x)\big)$.

\noindent\textit{Proof.} Pick $\varepsilon >0$. Since $H^\ast$ is continuous at $H(x)$, we can fix $\delta >0$ such that
$$
\di\Big(H^\ast\big(H(x)\big),H^\ast(y)\Big)<\frac{\varepsilon}{2}\text{ whenever }y\in X\text{ and }\di\big(H(x),y\big)<\delta.
$$
Since $(H_n:n\in\omega)$ converges to $H$, we can fix $N\in\omega$ such that
$$
\di\big(H_n(x),H(x)\big)<\delta\text{ whenever }n\geq N.
$$
Since $(H^{-1}_m:m\in\omega)$ converges uniformly to $H^\ast$, we can fix $M\in\omega$ such that
$$
\di\big(H^{-1}_m(y),H^\ast(y)\big)<\frac{\varepsilon}{2}\text{ whenever }m\geq M\text{ and }y\in X.
$$
Finally, set $K=\maxi\{M,N\}$. Then
\begin{align}
\di\Big(x_{m,n},H^\ast\big(H(x)\big)\Big)&=\di\Big(H_m^{-1}\big(H_n(x)\big),H^\ast\big(H(x)\big)\Big)\nonumber\\
&\leq\di\Big(H_m^{-1}\big(H_n(x)\big),H^\ast\big(H_n(x)\big)\Big)+\di\Big(H^\ast\big(H_n(x)\big),H^\ast\big(H(x)\big)\Big)\nonumber\\
&<\frac{\varepsilon}{2}+\frac{\varepsilon}{2}\nonumber\\
&=\varepsilon\nonumber
\end{align}
whenever $m,n\geq K$. $\blacksquare$

\noindent\textbf{Claim 3.} $y_{m,n}\to x$.

\noindent\textit{Proof.} Proceed as in the proof of Claim 1, except that condition $(1)$ should be used instead of condition $(2)$. $\blacksquare$

\noindent\textbf{Claim 4.} $y_{m,n}\to H\big(H^\ast(x)\big)$.

\noindent\textit{Proof.} Proceed as in the proof of Claim 2. $\blacksquare$
\end{proof}

\end{document}
